\section{组合数学基础}

\subsection{排列组合基础}

需要 \textbf{\texttt{ModInt}} 模板。

组合计数取模类，\texttt{MOD} 需为素数：\texttt{C(n,m)} 返回 $\binom{n}{m}\bmod MOD$，$n,m$ 可任意大（自动 Lucas 分解）。

\texttt{A(n,m)}、\texttt{fact(n)}、\texttt{inv(n)} 分别为排列数 $n!/(n-m)!$、$n!$、$(n!)^{-1}$，一般要求 $n<MOD$（更大时返回模意义下结果）；首次查询按需预计算阶乘，单次 $\mathcal{O}(1)$。

\subsubsection{自适应扩容版本}

\begin{lstlisting}
template <int MOD>
struct Combination{
    using Z = ModInt<MOD>;
private:
    int siz = 0;
    std::vector<Z> _fact,_inv;
    void expend(int size){
        if (siz >= size) return;
        _fact.resize(size),_inv.resize(size);
        for (int i = siz;i < size;i++) _fact[i] = _fact[i-1] * i;
        _inv[size-1] = 1 / _fact[size-1];
        for (int i = size - 2;i >= siz;i--) _inv[i] = _inv[i+1] * (i+1);
        siz = size;
    }
    void check(int size){
        if (siz >= size) return;
        int n = siz;
        while(n < size) n <<= 1;
        n = std::min<unsigned int>(n,MOD);
        expend(n);
    }
public:
    Combination():siz(1),_fact(1,1),_inv(1,1){}
    Z fact(int n){
        if (n < 0) return inv(-n);
        if (n >= MOD) return 0;
        check(n+1);
        return _fact[n];
    }
    Z inv(int n){
        if (n < 0) return fact(-n);
        if (n >= MOD) return 0;
        check(n+1);
        return _inv[n];
    }
    Z A(int n,int m){
        if (n < 0 || m < 0 || n < m) return 0;
        if (m >= MOD) return 0;
        if (n < MOD) {
            check(n+1);
            return _fact[n]*_inv[n-m];
        }
        check(m+1);
        return C(n,m)*_fact[m];
    }
    Z C(int n,int m){
        if (n < 0 || m < 0 || n < m) return 0;
        if (n < MOD && m < MOD) {
            check(n+1);
            return _fact[n] * _inv[m] * _inv[n - m];
        } else return C(n/MOD,m/MOD)*C(n%MOD,m%MOD);
    }
};
\end{lstlisting}

\subsubsection{赛时竞赛版本}

\begin{lstlisting}
template <int MOD>
struct Combination{
    using Z = ModInt<MOD>;
private:
    int siz = 0;
    constexpr static int N = 1e6+7;
    std::vector<Z> _fact,_inv;
public:
    Combination():siz(N),_fact(N,1),_inv(N,1){
        for (int i = 1;i < N;i++) _fact[i] = _fact[i-1] * i;
        _inv[N-1] = 1 / _fact[N-1];
        for (int i = N - 2;i >= 0;i--) _inv[i] = _inv[i+1] * (i+1);
    }
    Z fact(int n){
        if (n < 0) return inv(-n);
        return _fact[n];
    }
    Z inv(int n){
        if (n < 0) return fact(-n);
        return _inv[n];
    }
    Z A(int n,int m){
        if (n < 0 || m < 0 || n < m) return 0;
        return _fact[n]*_inv[n-m];
    }
    Z C(int n,int m){
        if (n < 0 || m < 0 || n < m) return 0;
        if (n < MOD && m < MOD) {
            return _fact[n] * _inv[m] * _inv[n - m];
        } else return C(n/MOD,m/MOD)*C(n%MOD,m%MOD);
    }
    Z Catalan(int n){
        if (n < 0) return 0;
        if (n == 0) return 1;
        return C(2*n,n) - C(2*n,n-1);
    }
};
\end{lstlisting}

\subsection{卡特兰数}

通常情况下有四种形式：

\begin{enumerate}
    \item $$C_n = \begin{cases}
        1 & ,n=0\\
        \sum_{i=0}^{n-1} C_iC_{n-i-1}&,n > 0\\
    \end{cases}$$
    \item $$C_n=\frac{1}{n+1}\binom{2n}{n}=\frac{(2n)!}{n!(n+1)!}$$
    \item $$C_n=\binom{2n}{n} - \binom{2n}{n+1}$$
    \item $$C_n=\frac{4n-2}{n+1}C_{n-1},n > 0,C_0 = 1$$
\end{enumerate}

代码为片段：\texttt{Z} 为全局取模类型，\texttt{C(n,m)} 为组合数查询函数（如 \texttt{Combination} 对象的 \texttt{C} 成员）。

\begin{lstlisting}
Z Catalan(int n){
    if (n < 0) return 0;
    if (n == 0) return 1;
    return C(2*n,n) - C(2*n,n-1);
}
\end{lstlisting}

\subsection{卢卡斯定理}

对于素数 $p$，有 $\binom{n}{k}\equiv \binom{\lfloor n/p\rfloor}{\lfloor k/p\rfloor}\binom{n\bmod p}{k\bmod p}\pmod p$。其中，当 $n<k$ 时，二项式系数 $\dbinom{n}{k}$ 规定为 $0$。
\subsection{扩展卢卡斯定理（exLucas）}

模数 $MOD$ 不是素数时 $\binom nm\bmod MOD$ 不能直接算（分母可能与 $MOD$ 不互素）。做法是把 $MOD$ 分解为质幂之积 $\prod p_i^{e_i}$，分别在模 $p_i^{e_i}$ 下求出 $\binom nm$，再用中国剩余定理合并（见数论章节）。

求 $\binom nm\bmod p^k$ 时把阶乘中的 $p$ 因子全部提出：由勒让德公式 $v_p(x!)=\sum_{i\ge1}\lfloor x/p^i\rfloor$ 得到分子、分母中 $p$ 的净次数；提出 $p$ 之后的 $x!$ 是一个长为 $p^k$ 的周期连乘积与 $\lfloor x/p\rfloor!$ 的乘积，于是可以递归求解。去掉 $p$ 因子后分母与 $p^k$ 互素，用扩展欧几里得求逆元消去。

\begin{lstlisting}
using i64 = long long;
using i128 = __int128;
namespace ExLucas {
    i64 exgcd(i64 a,i64 b,i64& x,i64& y){
        if (!b) return x = 1,y = 0,a;
        i64 g = exgcd(b,a % b,y,x);
        return y -= a / b * x,g;
    }
    i64 Pow(i64 a,i64 b,i64 mod){
        i64 res = 1;
        a %= mod;
        for (;b;b >>= 1,a = (i128)a * a % mod)
            if (b & 1) res = (i128)res * a % mod;
        return res;
    }
    i64 inv(i64 a,i64 mod){
        i64 x,y;
        exgcd(a,mod,x,y);
        return (x % mod + mod) % mod;
    }
    i64 fac(i64 n,i64 p,i64 pk){
        if (!n) return 1;
        i64 res = 1;
        for (i64 i = 1;i <= pk;i++)
            if (i % p) res = (i128)res * i % pk;
        res = Pow(res,n / pk,pk);
        for (i64 i = 1;i <= n % pk;i++)
            if (i % p) res = (i128)res * i % pk;
        return (i128)res * fac(n / p,p,pk) % pk;
    }
    i64 cnt(i64 n,i64 p){
        i64 res = 0;
        for (;n;n /= p) res += n;
        return res;
    }
    i64 C(i64 n,i64 m,i64 p,i64 pk){
        i64 e = cnt(n,p) - cnt(m,p) - cnt(n - m,p);
        i64 a = fac(n,p,pk),b = fac(m,p,pk),c = fac(n - m,p,pk);
        return (i128)Pow(p,e,pk) * a % pk * inv(b,pk) % pk * inv(c,pk) % pk;
    }
    i64 C(i64 n,i64 m,i64 mod){
        if (m < 0 || m > n) return 0;
        i64 M = mod,ans = 0;
        for (i64 p = 2;p * p <= mod;p++) if (mod % p == 0){
            i64 pk = 1;
            while (mod % p == 0) mod /= p,pk *= p;
            ans = (ans + (i128)C(n,m,p,pk) * (M / pk) % M * inv(M / pk,pk)) % M;
        }
        if (mod > 1) ans = (ans + (i128)C(n,m,mod,mod) * (M / mod) % M * inv(M / mod,mod)) % M;
        return ans;
    }
}
\end{lstlisting}

\texttt{ExLucas::C(n,m,MOD)} 返回 $\binom nm\bmod MOD$，$n,m$ 可达 $10^{18}$，$MOD$ 为 $10^{18}$ 以内的任意正整数；$m<0$ 或 $m>n$ 返回 $0$，$MOD=1$ 时恒返回 $0$。四参数重载 \texttt{C(n,m,p,pk)} 只算模质幂 $p^k$ 的结果，需要单独用时可直接调用。

复杂度为 $\mathcal{O}\big(\sum_i p_i^{e_i}\log_{p_i}n\big)$，瓶颈是模 $p_i^{e_i}$ 的预处理：$MOD$ 的质幂因子在一百万以内时很快；若 $MOD$ 本身是大素数（此时 $\sum p_i^{e_i}=MOD$）会退化成 $\mathcal{O}(MOD)$ 而过慢，这种情况应直接用素数模下的卢卡斯定理。
